\documentclass[../../main.tex]{subfiles}

\begin{document}
\begin{CJK*}{UTF8}{gbsn}

\section*{Exercise 4}

For a data set $\mathcal{D} = ((X_1,Y_1), \ldots, (X_m, Y_m)) \in (\mathcal{X} \times \mathcal{Y})^m$,
where $\mathcal{Y} = \mathbb{R}$, and any permutation $\sigma \in S_m$, put 

\begin{equation*}
    \mathcal{D}_{\sigma} = ((X_{\sigma(1)},Y_{\sigma(1)}), \ldots, (X_{\sigma(m)}, Y_{\sigma(m)})).
\end{equation*}
Given model $\hat{f}(x; \mathcal{D})$,
we construct a randomized score function $s((x,y); \mathcal{D}, \xi)$,
where $\xi$ is a random variable independent of the data set,
as follows.
Let $\sigma_{\xi}$ be an $S_m$-valued random variable 
depending only on $\xi$ that has the uniform distribution on $S_m$
if $\xi$ has the uniform distribution on $[0,1]$.
Put 

\begin{equation*}
    s((x,y); \mathcal{D}, \xi) = \abs{y - \hat{f}(x; \mathcal{D}_{\sigma_{\xi}})}
\end{equation*}
for all $(x,y)$ and all values of $\xi$.
Prove that $s$ is a symmetric score function in the following sense:
for every fixed test points $(x_1^{'}, y_1^{'}), \cdots, (x_{m'}^{'}, y_{m'}^{'})$
and every $\tau \in S_m$,
the vector 

\begin{equation*}
    \bracket{s\bracket{(x_1^{'}, y_1^{'}); \mathcal{D}, \xi}, \ldots, s\bracket{(x_{m'}^{'}, y_{m'}^{'}); \mathcal{D}, \xi}} 
    \in \mathbb{R}^{m'}
\end{equation*}
has the same distribution as 
\begin{equation*}
    \bracket{s\bracket{(x_1^{'}, y_1^{'}); \mathcal{D}_{\tau}, \xi}, \ldots, 
    s\bracket{(x_{m'}^{'}, y_{m'}^{'}); \mathcal{D}_{\tau}, \xi}} 
    \in \mathbb{R}^{m'}
\end{equation*}
when $\xi$ has the uniform distribution on $[0,1]$.

\paragraph{Solution}
For each $i \in \{1, \ldots, m'\}$,
the random variable 

\begin{equation*}
    s\bracket{(x_i^{'}, y_i^{'}); \mathcal{D}, \xi}
    = \abs{y_i^{'} - \hat{f}(x_i^{'}; \mathcal{D}_{\sigma_{\xi}})}
\end{equation*}
has the same distribution as 

\begin{equation*}
    \abs{y_i^{'} - \hat{f}(x_i^{'}; \mathcal{D}_{\tau \circ \sigma_{\xi}})}
    = s\bracket{(x_i^{'}, y_i^{'}); \mathcal{D}_{\tau}, \xi}
\end{equation*}
as long as $\xi$ has the uniform distribution on $[0,1]$
because, by construction, both $\sigma_{\xi}$ and $\tau \circ \sigma_{\xi}$
have uniform distribution on $S_m$.
////

\end{CJK*}
\end{document}
